\documentstyle[% 


righttag% 


,11pt% 


,amssymb% 


,russian% 


,izvru% 


]{amsart} 


 


%       Math definitions 


 


\theoremstyle{plain}


\theorembodyfont{\it}


\newtheorem{thmnonum}{\hskip\parindent Теорема}


\renewcommand{\thethmnonum}{}


 


\setcounter{page}{74} 


\god{1997} 


\nomer{7~(422)} 


\udk{517.958} 


\author{\it И.Н.\,СИДОРОВ} 


\title{ПОСТРОЕНИЕ ИНТЕГРАЛЬНОГО ПРЕДСТАВЛЕНИЯ \\ 


РЕШЕНИЯ УРАВНЕНИЙ РАВНОВЕСИЯ РАЗЛИЧНЫХ ВАРИАНТОВ \\ 


ТЕОРИИ ОБОЛОЧЕК СЛОЖНОЙ ГЕОМЕТРИИ} 


 


%\thanks{Работа выполнена при поддержке РФФИ, грант \No. 94-01-00183.} 


\date{} 


\pagestyle{plain} 


\begin{document} 


 


\maketitle 


 


В данной работе рассматривается изотропная тонкая или пологая оболочка 


постоянной толщины $2h$ в трехмерном евклидовом пространстве $R^3$. Оболочка 


имеет срединную поверхность $S$, удовлетворяющую необходимым требованиям 


гладкости, кусочно-гладкий контур этой поверхности $\Gamma$ и боковую 


поверхность $\Sigma$, образованную движением нормали $\bold m$ к поверхности 


$S$ вдоль контура $\Gamma$. В работе \cite{Vek} для такой оболочки получена 


конечная система уравнений равновесия в векторной форме в криволинейной 


нормальной системе координат, связанной с ее срединной поверхностью. Следуя 


обозначениям этой работы, представим конечную систему уравнений равновесия в 


виде 


\begin{equation} 


\frac{1}{\sqrt \alpha}


\partial_{\alpha,x}(\sqrt\alpha {\bold P}^{(k)\alpha}({\bold r}_x)) 


-\frac{1}{h}\underline{{\bold P}^{(k)3}}({\bold r}_x)+{\boldsymbol \Phi}^{(k)} 


({\bold r}_x) = 0,\quad {\bold r}_x\in S,\quad k=\overline{0,N}, 


\end{equation} 


где ${\bold r}_x$ --- радиус-вектор точек срединной поверхности $S$ 


\begin{gather*} 


{\bold P}^{(k)i}=\frac{2k+1}{2h}\int^h_{-h}P_k\left(\frac{z}{h}\right) 


{\bold P}^i dz,\\ 


\underline{{\bold P}^{(k)3}}=(2k+1)\sum^{[(k-1)/2]}_{i=0}{\bold P}^{(k-2i-1)3}, 


\quad k=\overline{1,N}, \quad \underline{{\bold P}^{(0)3}}=0,\\ 


{\boldsymbol \Phi^{(k)}}={\bold F}^{(k)}+[{\bold p}^{(+)}_n+(-1)^k 


{\bold p}^{(-)}_n]\frac{(2k+1)}{2h}, 


\end{gather*} 


${\bold P}^i$ --- векторные составляющие тензора напряжений, ${\bold F}$ 


--- вектор объемных сил, ${\bold p}^{(\pm)}_n$ --- 


заданный вектор напряжений на верхней (нижней) лицевой поверхности оболочки, 


$P_k\left(\frac{z}{h}\right )$ --- полином Лежандра $k$-го порядка, $z\in 


[-h,h]$ --- поперечная координата нормальной 


системы координат, $\partial_{\alpha,x}\equiv\frac{\partial}{\partial 


x^\alpha}$ $(x^1,$ $x^2$ --- криволинейные координаты), $a$ --- определитель 


метрического тензора срединной поверхности. В системе уравнений (1) и в 


последующих формулах греческие индексы пробегают значения 1, 2, латинские 


индексы --- 1, 2, 3, по повторяющимся индексам проводится 


суммирование (исключением 


является индекс, определяющий порядковый номер в разложении вектор-функции по 


полиномам Лежандра). 


 


В рамках $(N+1)$-моментного приближения вектор перемещений для элементов 


оболочки определяется формулой \cite{Vek} 


\begin{equation*} 


{\bold u}_N=\sum^N_{k=0}P_k\left(\frac{z}{h}\right ){\bold w}^{(k)}. 


\end{equation*} 


Тогда согласно обобщенному закону Гука 


\begin{equation} 


{\bold P}^{(k)i}=\lambda({\bold r}^m, D_{m,x}{\bold w}^{(k)}) 


{\bold r}^i +\mu({\bold r}^i, D_{m,x}{\bold w}^{(k)}){\bold r}^m+ 


\mu a^{im}(D_{m,x}{\bold w}^{(k)},{\bold r}^j ){\bold r}_j , 


\end{equation} 


где 


$$ 


D_{m,x}{\bold w}^{(k)}= 


\begin{cases} 


\partial_{\alpha,x}{\bold w}^{(k)},& m=\alpha ,\\ 


\frac{2k+1}{2h}({\bold w}^{(k+1)}+{\bold w}^{(k+3)}+\dotsb),& m=3, 


\end{cases} 


$$ 


${\bold r}_j,$ ${\bold r}^i$ --- векторы основного и взаимного базисов 


срединной поверхности оболочки, $\lambda,$ $\mu$ --- параметры Ламе, $a^{im}$ 


--- контравариантные компоненты метрического тензора. 


 


Для построения интегральной формы решения системы (1) воспользуемся 


уравнением, которому удовлетворяют моменты тензора перемещений Кельвина  


(\cite{Ugo}, c.\,43) $U^p_{(i)}({\bold R},{\boldsymbol \Lambda})$ $({\bold R} 


={\bold r}_x+z{\bold m}({\bold r}_x))$ 


\begin{gather} 


\partial_{\alpha,x}(\sqrt\alpha{\bold T}^{(k)\alpha}_{(i)} 


({\bold r}_x,{\boldsymbol \Lambda}))-\frac{1}{h}\sqrt\alpha 


\underline{{\bold T}^{(k)3}_{(i)}} ({\bold r}_x,{\boldsymbol \Lambda})+\\ 


%


+\frac{2k+1}{2h}\sqrt\alpha[{\bold e}_i P_k \left(\frac{\xi}{h}\right )


\delta({\bold r}_x-{\bold r}_\eta) + ({\bold T}^{(+)3}_{(i)} 


-(-1)^k {\bold T}^{(-)3}_{(i)})]=0,\notag\\ 


%


{\boldsymbol \Lambda} ={\bold r}_\eta+\xi{\bold m}({\bold r}_\eta), 


\quad {\bold r}_\eta\in S,\quad |\xi|<h, \notag


\end{gather} 


где $\delta({\bold r}_x-{\bold r}_\eta)$ --- дельта-функция Дирака, 


${\bold e}_i$ --- единичный вектор 


декартовой системы 


координат, определяющий направление действия сосредоточенной силы в 


бесконечной упругой среде и приложенный в точке с радиус-вектором 


${\boldsymbol \Lambda}$,


$ {\bold T}^{(\pm)3}_{(i)}={\bold T}^3_{(i)}({\bold r}_x\pm h{\bold m} 


({\bold r}_x), {\boldsymbol \Lambda}),$ 


\begin{equation} 


{\bold T}^{(k)j}_{(i)}=\lambda({\bold r}^m, D_{m,x} 


{\bold U}^{(k)}_{(i)}){\bold r}^j +\mu({\bold r}^j, D_{n,x} 


{\bold U}^{(k)}_{(i)}){\bold r}^n 


+\mu a^{jm}(D_{m,x}{\bold U}^{(k)}_{(i)},{\bold r}^p){\bold r}_p, 


\quad {\bold U}^{(k)}_{(i)}=U^{(k)p}_{(i)}{\bold e}_p. 


\end{equation} 


Система уравнений (3) получена с учетом предположения, что для тонких или 


пологих оболочек с погрешностью $zb^\beta_\alpha$ по сравнению 


с $\delta^\beta_\alpha=({\bold r}^\beta,{\bold r}_\alpha)$ 


$(b^\beta_\alpha$ --- смешанные компоненты 


второго метрического тензора срединной поверхности оболочки) векторы базиса и 


определитель метрического тензора элементов оболочки можно заменить на 


соответствующие величины срединной поверхности. Заключением индекса $i$ в 


круглые скобки подчеркивается, что вектор Кельвина ${\bold U}_{(i)}$ 


соответствует единичной 


силе ${\bold e}_i$ в декартовой системе координат, 


а уравнение для этого вектора 


записывается в криволинейной системе координат. 


 


Предположим, что вектор ${\bold u}_N$ 


является решением системы (1). Тогда интеграл 


\begin{align*} 


I &=\int^h_{-h}dz \iint_S\biggl(\sum^N_{k=0}\partial_{\alpha,x} 


(\sqrt\alpha{\bold P}^{(k)\alpha})P_k\left(\frac{z}{h}\right ), 


{\bold U}_{(i)}\biggr)dx^1 dx^2 =\\ 


%


&=\iint_S\sum^N_{k=0}\frac{2h}{2k+1}(\partial_{\alpha,x} 


({\bold P}^{(k),\alpha}),{\bold U}^{(k)}_{(i)})dx^1 dx^2, 


\end{align*} 


который в соответствии с (1) равен 


\begin{equation} 


I=\iint_S\sum^N_{k=0}\frac{2h}{2k+1}\left(\frac{1}{h}


\underline{{\bold P}^{(k)3}}- 


{\boldsymbol \Phi}^{(k)}, {\bold U}^{(k)}_{(i)}\right)\sqrt a\,dx^1 dx^2, 


\end{equation} 


преобразуем к виду 


\begin{equation} 


I=\int_\Gamma\sum^N_{k=0}\frac{2h}{2k+1}({\bold P}^{(k)\alpha}, 


{\bold U}^{(k)}_{(i)})n^\Sigma_\alpha ds_t -


\iint_S\sqrt a\sum^N_{k=0}\frac{2h}{2k+1} 


({\bold P}^{(k)\alpha},\partial_{\alpha,x}{\bold U}^{(k)}_{(i)})


dx^1 dx^2, 


\end{equation} 


где $ds_t$ --- элемент дуги контура $\Gamma$ срединной поверхности, 


$n^\Sigma_\alpha$ --- компоненты 


единичной нормали к боковой поверхности оболочки $\Sigma$. 


 


Сумму скалярных произведений во втором интеграле правой части (6) с 


учетом (2) и (4) представим следующим образом: 


\begin{multline} 


\sum^N_{k=0}\frac{2h}{2k+1}


({\bold P}^{(k)\alpha},\partial_{\alpha,x}{\bold U}^{(k)}_{(i)})=


\sum^N_{k=0}\frac{2h}{2k+1}


({\bold T}^{(k)\alpha}_{(i)},\partial_{\alpha,x}{\bold w}^{(k)})+\\


+\sum^N_{k=0}\frac{2h}{2k+1}


[({\bold T}^{(k)3}_{(i)},D_{3,x}{\bold w}^{(k)})-


({\bold P}^{(k)3},D_{3,x}{\bold U}^{(k)}_{(i)})]. 


\end{multline} 


Используя (7), соотношение 


\begin{multline*} 


\sum^N_{k-0}\frac{2h}{2k+1}\sqrt a 


({\bold T}^{(k)\alpha}_{(i)},\partial_{\alpha,x}{\bold w}^{(k)}) 


=\sum^N_{k=0}\frac{2h}{2k+1}\partial_{\alpha,x}(\sqrt a


({\bold T}^{(k)\alpha}_{(i)},{\bold w}^{(k)}))-\\ 


%


-\sum^N_{k=0}\frac{2h}{2k+1}(\partial_{\alpha,x}(\sqrt a


({\bold T}^{(k)\alpha}_{(i)}),{\bold w}^{(k)}),


\end{multline*} 


а также (3), правую часть в (6) преобразуем к виду 


\begin{multline} 


I=\sum^N_{k=0}\frac{2h}{2k+1}\biggl\{\int_\Gamma[ 


({\bold P}^{(k)\alpha},{\bold U}^{(k)}_{(i)})- 


({\bold T}^{(k)\alpha}_{(i)},{\bold w}^{(k)})]n^\Sigma_\alpha ds_t- \\


%


-\iint_S\sqrt a[({\bold T}^{(k)3}_{(i)},D_{3,x}{\bold w}^{(k)})- 


({\bold P}^{(k)3},D_{3,x}{\bold U}^{(k)}_{(i)})+


\frac{2k+1}{2h}P_k(\xi/h)({\bold e}_i \delta({\bold r}_x-{\bold r}_\eta), 


{\bold w}^{(k)})-\\


%


-\frac{1}{h}({\bold T}^{(k)3}_{(i)},{\bold w}^{(k)})+


\frac{2k+1}{2h}({\bold T}^{(+)3}_{(i)}-(-1)^k 


{\bold T}^{(-)3}_{(i)},{\bold w}^{(k)})]dx^1 dx^2\biggr\}. 


\end{multline} 


После некоторых преобразований правой части (8) с учетом (5) получим 


интегральное представление вектора перемещений ${\bold u}_N$


\begin{gather} 


\beta_{ip}({\bold r}_\eta)({\bold e}_p,{\bold u}_N({\bold r}_\eta))= 


\sum^N_{k=0}\biggl\{\iint_S\biggl[\frac{2h}{2k+1} 


({\boldsymbol \Phi}^{(k)},{\bold U}^{(k)}_{(i)})- 


({\bold T}^{(+)3}_{(i)}-(-1)^k{\bold T}^{(-)3}_{(i)},{\bold w}^{(k)})+ 


({\bold P}^{(k)3},{\bold U}^{(+)}_{(i)}- 


\notag \\ 


-(-1)^k{\bold U}^{(-)}_{(i)} 


-\sum^N_{j=0}{\bold U}^{(j)}_{(i)}[1-(-1)^{k+j}])\biggr]dS_R+


\int_\Gamma\frac{2h}{2k+1}[({\bold P}^{(k)\Sigma}_n, 


{\bold U}^{(k)}_{(i)})-({\bold P}^{(k)}_{(i)},{\bold w}^{(k)})] 


ds_t\biggr\}, 


\end{gather} 


$$ 


\beta_{ip}(r_\eta)= 


\begin{cases} 


\delta_{ip},   & {\bold r}_\eta\in S; \\


          0,   & {\bold r}_\eta\notin S,\text{\ }  |\xi|<h; \\


1/2\delta_{ip},& {\bold r}_\eta\in \Gamma, 


\end{cases} 


$$ 


где ${\bold P}^{(k)\Sigma}_n={\bold P}^{(k)\alpha}n^\Sigma_\alpha,$ 


${\bold P}^{(k)}_{(i)}={\bold T}^{(k)\alpha}_{(i)}n^\Sigma_\alpha 


=P^{(k)p}_{(i)}{\bold e}_p$ $(P^p_{(i)}$ --- тензор напряжений Кельвина 


\cite{Ugo}, с.\,43),  


${\bold U}^{(\pm)}_{(i)}={\bold U}_{(i)}({\bold r}_x \pm h{\bold m} 


({\bold r}_x),{\boldsymbol \Lambda}),$ $\delta_{ip}$ 


--- символ Кронекера. При выводе (9) были использованы соотношения 


\begin{multline*} 


\sum^N_{k=0}\frac{2h}{2k+1} 


[({\bold P}^{(k)3},D_{3,x}{\bold U}^{(k)}_{(i)})-\frac{1}{h}


({\bold P}^{(k)3},{\bold U}^{(k)}_{(i)})]= \\ 


%


=\sum^N_{k=0}({\bold P}^{(k)3},{\bold U}^{(+)}_{(i)})- 


(-1)^k{\bold U}^{(-)}_{(i)})- 


\sum^N_{k=0}\,\sum^N_{j=0} 


({\bold P}^{(k)3},{\bold U}^{(j)}_{(i)})[1-(-1)^{k+j}], 


\end{multline*} 


%


\begin{multline*}


\sum^N_{k=0}\frac{2h}{2k+1}({\bold T}_{(i)}^{(k)3},D_{3,x}


{\bold w^{(k)}})=\int^h_{-h}\Bigl({\bold T^3_{(i)}},


\sum^N_{k=0}P_k D_{3,x}{\bold w^{(k)}}\Bigr)dz= \\


%


=\int^h_{-h}\biggl({\bold T^3_{(i)}},\frac{\partial}{\partial z}


\Bigl(\,\sum^N_{k=0}P_k {\bold w^{(k)}}\Bigr)\biggr)dz=


\sum^N_{k=0}\frac{2}{2k+1}({\bold T}_{(i)}^{(k)3},{\bold w^{(k)}}).


\end{multline*}


 


Покажем, что вектор перемещений ${\bold u}_N,$ удовлетворяющий системе 


соотношений (9), является решением системы (1). Для этого подействуем на обе 


части этих соотношений оператором 


\begin{gather} 


\displaybreak[3]


L^{(N)}_\eta{\bold u}_N=\sum^N_{k=0}P_k\left(\frac{\xi}{h}\right) 


L^{(k)}_\eta {\bold u}_N,\\ 


%


L^{(k)}_\eta{\bold u}_N=\frac{1}{\sqrt a}\partial_{\alpha,\eta} 


\bigl\{\sqrt a [\lambda({\bold r}^m, D_{m,\eta}{\bold w}^{(k)} 


({\bold r}_\eta)){\bold r}^\alpha+


\mu({\bold r}^\alpha,D_{n,\eta}{\bold w}^{(k)}({\bold r}_\eta)) 


{\bold r}^n+\mu a^{\alpha m}(D_{m, \eta}{\bold w}^{(k)}({\bold r}_\eta), 


{\bold r}^j){\bold r}_j]\bigr\}-\notag\\ 


%


-\frac{1}{h}\underline{{\bold P}^{(k)3}_N}({\bold r}_\eta)+\frac{2k+1}{2h} 


[{\bold P}^{(+)3}_N({\bold r}_\eta)-(-1)^k {\bold P}^{(-)3}_N 


({\bold r}_\eta)],\notag\\ 


%


{\bold P}^j_N=\sum^N_{k=0}P_k\left(\frac{\xi}{h}\right) 


{\bold P}^{(k)j}_N\quad ({\bold r}_\eta\in S). \notag


\end{gather} 


Принимая во внимание, что имеют место равенства 


\begin{align} 


&\underline{{\bold P}^{(k)j}_N}({\bold u}_N)=


\underline{{\bold P}^{(k)j}}({\bold u}_N),\notag\\ 


%


& {\bold P}^{(+)3}_N-(-1)^k{\bold P}^{(-)3}_N=\sum^N_{j=0} 


[1-(-1)^{k+j}]{\bold P}^{(j)3},\\ 


%


& L^{(k)}_\eta{\bold U}^{(m)}_{(i)}({\bold r}_x,{\boldsymbol \Lambda}) 


=-\frac{2m+1}{2h}\delta_{km}{\bold e}_i\delta({\bold r}_x-{\bold r}_\eta), 


\quad m=\overline{0,N},\notag\\ 


%


& L^{(k)}_\eta{\bold U}^{(m)}_{(i)}({\bold r}_{x,\Gamma},{\boldsymbol \Lambda}) 


=L^{(k)}_\eta{\bold P}^{(m)}_{(i)}({\bold r}_{x,\Gamma},{\boldsymbol \Lambda})=0, 


\notag\\ 


%


& L^{(k)}_\eta[{\bold T}^{(+)3}_{(i)}-(-1)^k 


{\bold T}^{(-)3}_{(i)}]= 


L^{(k)}_\eta[{\bold U}^{(+)}_{(i)}-(-1)^k 


{\bold U}^{(-)}_{(i)}]=0,\notag 


\end{align} 


$$


\sum^N_{k=0}\sum^N_{j=0}[1-(-1)^{k+j}]({\bold P}^{(j)3}, 


{\bold U}^{(k)}_{(i)})= 


\sum^N_{k=0}\sum^N_{j=0}[1-(-1)^{k+j}]({\bold P}^{(k)3}, 


{\bold U}^{(j)}_{(i)}),


$$


после действия оператора $L^{(N)}_\eta$ на интегральное представление вектора 


${\bold u}_N,$ из (10) и (11) имеем 


$$ 


\sum^N_{k=0}\left[\frac{1}{\sqrt a}\partial_{\alpha,\eta} 


(\sqrt a {\bold P}^{(k)\alpha}({\bold r}_\eta))- 


\frac{1}{h}\underline{{\bold P}^{(k)3}}({\bold r}_\eta)+{\boldsymbol \Phi}^{(k)} 


({\bold r}_\eta)\right]P_k\left(\frac{\xi}{h}\right)=0. 


$$ 


Из последнего равенства в силу независимости полиномов Лежандра следует (1). 


 


Таким образом, доказана 


 


\begin{thmnonum} Для того чтобы вектор ${\bold u}_N$ являлся решением 


системы {\em(1),} необходимо и достаточно, чтобы этот вектор удовлетворял 


системе соотношений {\em (9).} 


\end{thmnonum}


 


Следует отметить, что система соотношений (9) может быть использована при 


построении граничных интегральных уравнений для вычисления параметров 


напряженно-деформированного состояния элементов оболочки. В качестве способа 


получения этих уравнений можно предложить разложение (9) по полиномам 


Лежандра от поперечной координаты оболочки $\xi$. При этом получается система 


$N+1$ векторных интегральных уравнений относительно $N+1$ неизвестных 


моментов вектора перемещений на срединной поверхности или вектора напряжений 


на контуре оболочки. 


 


\begin{thebibliography}{9} 


 


\bibitem{Vek} 


Векуа И.Н. {\em Некоторые общие методы построения различных вариантов 


теории оболочек}. -- М.: Наука, 1982. -- 286 с. 


\bibitem{Ugo} 


Угодчиков А.Г., Хуторянский Н.М. {\em Метод граничных элементов в механике 


деформируемого тела}. -- Казань: Изд-во Казанск. ун-та, 1986. -- 295 с. 


 


\end{thebibliography} 


 


\vskip 0.5cm 


 


\post{\it Казанский государственный}{\it Поступила} 


\post{\it технический университет}{30.01.1996} 


 


\end{document} 


